\documentclass[12pt,a4paper]{article}

\usepackage{cmap}
\usepackage[T2A]{fontenc}
\usepackage[utf8]{inputenc}
\usepackage[russian]{babel}
\usepackage{amsmath,amssymb}
\usepackage[a4paper,margin=2cm]{geometry}

\newcommand{\CC}{\mathbb{C}}
\newcommand{\R}{\mathbb{R}}

\setlength{\parindent}{0pt}
\setlength{\parskip}{0.45em}
\sloppy

\begin{document}

\begin{center}
{\Large\bfseries Билет 12}
\end{center}

\noindent\fbox{%
  \begin{minipage}{\dimexpr\linewidth-2\fboxsep-2\fboxrule\relax}
  \normalfont
Степенные ряды с комплексными членами. Круг и радиус сходимости. Характер
сходимости степенного ряда в круге сходимости. Формула Коши--Адамара.
Сохранение радиуса сходимости степенного ряда при почленном дифференцировании
и интегрировании ряда. Первая теорема Абеля.
  \end{minipage}%
}

\section*{Краткий план}

Сначала степенной ряд с центром \(w_0\) сводится заменой \(z=w-w_0\) к ряду
\(\sum c_kz^k\). Радиус сходимости задается формулой Коши--Адамара
\[
  \frac1R=\limsup_{k\to\infty}\sqrt[k]{|c_k|}.
\]
Из обобщенного признака Коши следует характер сходимости: внутри круга ряд
сходится абсолютно, вне круга расходится, а на границе возможны оба случая.
На каждом замкнутом круге \(|z|\le r<R\) сходимость равномерна. Затем
доказывается первая теорема Абеля: из сходимости в точке \(z_0\) следует
абсолютная сходимость во всякой точке меньшего модуля. Наконец показывается,
что формальные почленное дифференцирование и интегрирование не меняют радиус
сходимости.

\section*{Определения}

\textbf{Степенной ряд.}
Пусть \(\{c_k\}_{k=0}^\infty\subset\CC\) и \(w_0\in\CC\). Ряд
\[
  \sum_{k=0}^{\infty}c_k(w-w_0)^k
\]
называется степенным рядом с комплексными членами. После замены \(z=w-w_0\)
его обычно записывают в виде
\[
  \sum_{k=0}^{\infty}c_kz^k.
\]

\textbf{Радиус и круг сходимости.}
Радиусом сходимости степенного ряда называется число
\[
  R\in[0,+\infty]\quad\text{такое, что}\quad
  \frac1R=\limsup_{k\to\infty}\sqrt[k]{|c_k|},
\]
с соглашениями \(1/0=+\infty\), \(1/(+\infty)=0\). Эта формула называется
формулой Коши--Адамара. Кругом сходимости называется множество
\[
  D_R=\{z\in\CC: |z|<R\}.
\]
Если \(R=+\infty\), то кругом сходимости считается вся комплексная плоскость
\(\CC\).

\textbf{Почленное дифференцирование и интегрирование.}
Формально из ряда \(\sum c_kz^k\) получаются ряды
\[
  \sum_{k=1}^{\infty}kc_kz^{k-1},
  \qquad
  \sum_{k=0}^{\infty}\frac{c_k}{k+1}z^{k+1}.
\]
Первый называется рядом, полученным почленным дифференцированием, второй --
рядом, полученным почленным интегрированием исходного степенного ряда.

\section*{Теоремы}

\textbf{1. Теорема о круге сходимости.}
Пусть \(R\) -- радиус сходимости ряда
\[
  \sum_{k=0}^{\infty}c_kz^k.
\]
Тогда:
\[
  |z|<R \quad\Longrightarrow\quad
  \sum_{k=0}^{\infty}c_kz^k \text{ сходится абсолютно},
\]
\[
  |z|>R \quad\Longrightarrow\quad
  \sum_{k=0}^{\infty}c_kz^k \text{ расходится}.
\]
На границе \(|z|=R\) общий вывод невозможен: в разных точках границы ряд может
сходиться или расходиться.

\textbf{2. Равномерная сходимость внутри круга.}
Если \(R>0\), то для любого \(r\in(0,R)\) ряд
\[
  \sum_{k=0}^{\infty}c_kz^k
\]
сходится равномерно на замкнутом круге \(\{z\in\CC: |z|\le r\}\).

\textbf{3. Первая теорема Абеля.}
Если степенной ряд
\[
  \sum_{k=0}^{\infty}c_kz^k
\]
сходится в точке \(z_0\), то в любой точке \(z_1\), для которой
\[
  |z_1|<|z_0|,
\]
этот ряд сходится абсолютно.

\textbf{4. Сохранение радиуса сходимости.}
Если \(R\) -- радиус сходимости ряда
\[
  \sum_{k=0}^{\infty}c_kz^k,
\]
то ряды
\[
  \sum_{k=1}^{\infty}kc_kz^{k-1},
  \qquad
  \sum_{k=0}^{\infty}\frac{c_k}{k+1}z^{k+1}
\]
имеют тот же радиус сходимости \(R\).

\section*{Доказательства}

\textbf{Теорема о круге сходимости.}
Пусть \(z\ne0\). Исследуем абсолютную сходимость ряда
\[
  \sum_{k=0}^{\infty}|c_kz^k|
\]
по обобщенному признаку Коши. Имеем
\[
  q=\limsup_{k\to\infty}\sqrt[k]{|c_kz^k|}
  =
  |z|\limsup_{k\to\infty}\sqrt[k]{|c_k|}
  =
  \frac{|z|}{R}.
\]
Если \(0<|z|<R\), то \(q<1\), поэтому ряд из модулей сходится, значит
исходный степенной ряд сходится абсолютно. При \(z=0\) ряд также сходится,
так как все члены начиная с первого равны нулю.

Если \(|z|>R\), то \(q>1\). По признаку Коши члены ряда
\(|c_kz^k|\) не стремятся к нулю. Следовательно, и \(c_kz^k\) не стремятся к
нулю, поэтому необходимое условие сходимости не выполнено и ряд расходится.

На границе поведение не определяется одним только радиусом. Например,
\[
  \sum_{k=0}^{\infty}\frac{z^k}{k+1}
\]
имеет радиус сходимости \(R=1\). При \(z=1\) получается гармонический ряд
\(\sum_{k=1}^{\infty}1/k\), который расходится; при \(z=-1\) получается
знакочередующийся гармонический ряд, который сходится по признаку Лейбница.

\textbf{Равномерная сходимость внутри круга.}
Пусть \(0<r<R\) и \(|z|\le r\). Тогда
\[
  |c_kz^k|\le |c_k|r^k.
\]
Так как \(r<R\), числовой ряд
\[
  \sum_{k=0}^{\infty}|c_k|r^k
\]
сходится по уже доказанной теореме о круге сходимости. По признаку
Вейерштрасса ряд \(\sum c_kz^k\) сходится равномерно на \(|z|\le r\).

\textbf{Первая теорема Абеля.}
Пусть ряд \(\sum c_kz^k\) сходится при \(z=z_0\). Если \(z_0=0\), утверждение
пусто, потому что нет точки \(z_1\) с \(|z_1|<0\). Пусть \(z_0\ne0\).
Если бы \(R<|z_0|\), то точка \(z_0\) лежала бы вне круга сходимости, а по
теореме о круге сходимости ряд в ней расходился бы. Значит
\[
  R\ge |z_0|.
\]
Для всякого \(z_1\) с \(|z_1|<|z_0|\) получаем \(|z_1|<R\). Поэтому по
теореме о круге сходимости ряд в точке \(z_1\) сходится абсолютно.

\textbf{Сохранение радиуса при почленном дифференцировании.}
Сначала сравним радиусы рядов
\[
  \sum_{k=0}^{\infty}c_kz^k
  \quad\text{и}\quad
  \sum_{k=1}^{\infty}kc_kz^k.
\]
По формуле Коши--Адамара для второго ряда
\[
  \frac1{R_1}
  =
  \limsup_{k\to\infty}\sqrt[k]{|kc_k|}
  =
  \lim_{k\to\infty}\sqrt[k]{k}\,
  \limsup_{k\to\infty}\sqrt[k]{|c_k|}
  =
  \frac1R,
\]
поскольку \(\sqrt[k]{k}\to1\). Значит \(R_1=R\).

Ряды
\[
  \sum_{k=1}^{\infty}kc_kz^k
  \quad\text{и}\quad
  \sum_{k=1}^{\infty}kc_kz^{k-1}
\]
при \(z\ne0\) отличаются умножением суммы частичных сумм на постоянный
множитель \(1/z\), а при \(z=0\) оба вопроса сходимости решаются отдельно и
не меняют радиуса. Поэтому формально продифференцированный ряд имеет тот же
радиус \(R\).

\textbf{Сохранение радиуса при почленном интегрировании.}
Ряд
\[
  \sum_{k=0}^{\infty}\frac{c_k}{k+1}z^{k+1}
\]
после почленного дифференцирования дает исходный ряд
\(\sum c_kz^k\). По уже доказанному почленное дифференцирование не меняет
радиус сходимости. Следовательно, радиус интегрального ряда равен радиусу
исходного ряда, то есть \(R\).

\section*{Финальная устная версия}

Степенной ряд с комплексными членами имеет вид
\[
  \sum_{k=0}^{\infty}c_k(w-w_0)^k,\qquad c_k,w_0\in\CC.
\]
Заменой \(z=w-w_0\) он сводится к ряду \(\sum c_kz^k\). Его радиус
сходимости задается формулой Коши--Адамара:
\[
  \frac1R=\limsup_{k\to\infty}\sqrt[k]{|c_k|},
\]
где \(1/0=+\infty\), \(1/(+\infty)=0\). Круг сходимости -- это
\(\{z: |z|<R\}\), а при \(R=+\infty\) вся плоскость \(\CC\).

По обобщенному признаку Коши для ряда модулей
\[
  \limsup_{k\to\infty}\sqrt[k]{|c_kz^k|}=\frac{|z|}{R}.
\]
Поэтому при \(|z|<R\) ряд сходится абсолютно, а при \(|z|>R\) его члены не
стремятся к нулю и ряд расходится. На границе \(|z|=R\) возможны оба случая:
\(\sum z^k/(k+1)\) при \(z=1\) расходится, а при \(z=-1\) сходится. Кроме
того, на каждом замкнутом круге \(|z|\le r<R\) ряд сходится равномерно по
признаку Вейерштрасса, так как \(|c_kz^k|\le |c_k|r^k\).

Первая теорема Абеля говорит: если \(\sum c_kz^k\) сходится в точке \(z_0\),
то в любой точке \(z_1\) с \(|z_1|<|z_0|\) он сходится абсолютно.
Действительно, сходимость в \(z_0\) запрещает неравенство \(R<|z_0|\);
следовательно, \(R\ge |z_0|\), и точка \(z_1\) лежит строго внутри круга
сходимости.

Наконец, формальные ряды
\[
  \sum_{k=1}^{\infty}kc_kz^{k-1},
  \qquad
  \sum_{k=0}^{\infty}\frac{c_k}{k+1}z^{k+1}
\]
имеют тот же радиус сходимости. Для производного ряда это следует из
\[
  \limsup_{k\to\infty}\sqrt[k]{|kc_k|}
  =
  \lim_{k\to\infty}\sqrt[k]{k}\,
  \limsup_{k\to\infty}\sqrt[k]{|c_k|}
  =
  \limsup_{k\to\infty}\sqrt[k]{|c_k|}.
\]
Интегральный ряд имеет тот же радиус, потому что его почленное
дифференцирование возвращает исходный ряд, а дифференцирование радиус не
меняет.

\section*{Источники}

Иванов Г.Е. \emph{Лекции по математическому анализу. Часть 1}:
\texttt{IvGE\_1.pdf}, стр. 303--309.

Основные роли источника: стр. 303--304 -- определение комплексного
степенного ряда, радиус и круг сходимости, формула Коши--Адамара;
стр. 304--305 -- теорема о круге сходимости и первая теорема Абеля;
стр. 307 -- равномерная сходимость на замкнутых кругах внутри круга
сходимости; стр. 308--309 -- сохранение радиуса сходимости при почленном
дифференцировании и интегрировании.

Дополнительно был просмотрен \texttt{IvGE\_2.pdf}: прямого материала по
степенным рядам для этого билета там нет.

\end{document}
